% What Irregularity Costs.
%
% Build:  tectonic paper/main.tex
% Tables: paper/gen_tables.py results/sweep_YYYYmmdd_HHMMSS/sweep.csv
%
% No number in this paper is typed by hand. Every table under paper/tables/ is
% generated from the raw sweep CSVs, because five published figures in this
% project were later corrected and a hand-copied table would have preserved all
% of them.
\documentclass[11pt]{article}

\usepackage[margin=1in]{geometry}
\usepackage{booktabs}
\usepackage{amsmath}
\usepackage{graphicx}
\usepackage{listings}
\usepackage{xcolor}
\usepackage[hidelinks]{hyperref}
\usepackage{microtype}

\lstset{
  basicstyle=\ttfamily\footnotesize,
  breaklines=true,
  frame=single,
  framesep=4pt,
  rulecolor=\color{gray!40},
  columns=fullflexible,
}

\title{What Irregularity Costs:\\
CUDA C++, Rust, and Triton on a Hash-Blocked GPU Workload}

\author{Petr Korolev}
\date{\today}

\begin{document}
\maketitle

\begin{abstract}
GPU language comparisons are almost always run on tiled dense linear algebra,
where every toolchain is good and the differences are small. We implement the
same hash-blocked TSDF fusion kernel five times, in CUDA C++, in Rust through
NVIDIA's \texttt{cuda-oxide}, and in Triton, and measure it on a workload with
the opposite character: an open-addressed hash table with compare-exchange
insertion, data-dependent per-lane probe depth, and contended scatter.

The result is a split. On the regular stage, which walks a truncation band and
accumulates, all three languages land within a small factor of each other. On
the irregular stage, which probes and inserts, Rust stays close to hand-written
CUDA C++ while Triton is more than an order of magnitude slower. Language choice
is nearly free on the work that is usually benchmarked and expensive on the work
that is not.

We attribute both gaps to specific things the languages cannot express, rather
than reporting ratios. Triton's cost follows from a probe loop that must run to a
compile-time bound and from \texttt{tl.atomic\_cas} taking no mask, which forces
a scratch structure with no counterpart in the CUDA implementation; sizing that
structure badly also silently destroys the kernel's ability to use additional
streaming multiprocessors, a failure invisible on a single device. Rust's cost
was not visible in any instruction count: the Rust kernel issues fewer
instructions, fewer compare-exchanges, and fewer registers at identical
occupancy, and was still slower. Hardware counters located it in L1 residency. A
GPU-scope atomic load must be coherent across SMs, no NVIDIA L1 is coherent
across SMs, and so the type-correct way to read a shared location bypasses the
cache on every access. Replacing two such reads with plain loads, correct because
the algorithm rather than the memory model provides the guarantee, closes most of
the gap.

We report the measurement discipline as a result in its own right. Roughly a
dozen errors were caught during this work, three of which would have produced a
reversed or null conclusion rather than a wrong magnitude, and two of which were
caught only after the wrong number had been written down. We also report a
defect found and fixed in \texttt{cuda-oxide} itself: its scoped atomic load and
store could not be called at all in the build mode that produces real kernels.
\end{abstract}

\section{Introduction}

Pick up almost any comparison of GPU programming languages and the benchmark
will be a matrix multiply, a convolution, a stencil, or a reduction. These are
reasonable choices. They are also the workloads every toolchain has been tuned
for, and they share a property that makes them easy: the work per thread is
known before the kernel launches.

This paper measures the other case. A hash-blocked truncated signed distance
function (TSDF) volume, the data structure underneath most real-time 3D
reconstruction systems, is built by a kernel with none of that regularity. Each
point walks a band of voxels, hashes each one to a block coordinate, probes an
open-addressed table for a variable number of steps, and either finds an existing
block or claims one with a compare-exchange. Threads in the same warp do
different amounts of work, contend for the same table slots, and must observe one
another's insertions.

We implement this kernel five times. Two of the implementations are CPU and
third-party baselines used to establish correctness; three are the comparison
proper: CUDA C++, Rust compiled through NVIDIA's \texttt{cuda-oxide}, and
Triton. All three operate on identical device memory with an identical hash
layout, so the comparison measures how each language expresses the algorithm
rather than which algorithm each language made convenient.

\paragraph{Contributions.}
\begin{enumerate}
  \item A measured split between regular and irregular work
        (Section~\ref{sec:results}): the three languages are close on the stage
        that walks and accumulates, and separated by more than an order of
        magnitude on the stage that probes and inserts.
  \item Attribution of both gaps to language-level causes
        (Section~\ref{sec:attribution}), including one that no instruction count
        reveals and that required hardware performance counters to find.
  \item A catalogue of what each language cannot say
        (Section~\ref{sec:expressiveness}), each item encountered while
        implementing the algorithm rather than constructed to make a point.
  \item A measurement methodology and its failure modes
        (Section~\ref{sec:method}), reported because three of the errors it
        caught would have produced a confidently wrong conclusion rather than a
        noisy one.
\end{enumerate}

\paragraph{What this paper is not.} It is not a claim that one language is
better. Rust lands within a few percent of hand-written CUDA C++ on real data
once a single idiom is corrected, and Triton's cost is concentrated in one stage
that a practitioner could write in something else. It is a claim about where the
costs are, and why.

\section{The workload}
\label{sec:workload}

A truncated signed distance function (TSDF) volume stores, at each voxel, the
signed distance to the nearest observed surface, truncated to a band around it.
Fusing a depth image means walking that band along each point's view ray and
accumulating a weighted average. The surface is later recovered as the zero
crossing. It is the representation underneath KinectFusion and most real-time
reconstruction systems since.

A dense volume is unaffordable at useful resolution, so implementations store
only the blocks a surface actually touches, indexed by a spatial hash. That
choice is what makes the kernel irregular, and it is why this workload is the
vehicle for this paper rather than an application of it.

\subsection{Structure}

The volume is divided into blocks of $8^3 = 512$ voxels. A hash table maps block
coordinates to indices in a flat pool of voxel data; five parallel float arrays
hold the truncated distance, the accumulated weight, and three colour channels.
The table is open-addressed with linear probing and sized to the next power of
two above twice the pool capacity, so an occupied table is at most half full.

Each entry is 16 bytes: a 64-bit key and a 32-bit block index, with four bytes
of padding. The key packs the three block coordinates into one word, each axis
biased by $2^{20}$ into 21 bits:
%
\begin{equation*}
\mathrm{key}(x,y,z) = (x + 2^{20}) \ll 42 \;\vert\; (y + 2^{20}) \ll 21 \;\vert\; (z + 2^{20}).
\end{equation*}
%
Packing is not a space optimisation. It is what allows a single compare-exchange
to publish an entire coordinate. An earlier layout stored the three axes
separately and published them with a compare-exchange on the first followed by
plain stores to the other two; a reader could then match the first word,
mismatch on the rest, and probe onward, which at GPU thread counts degenerates
into a full-table scan per lookup. The 64-bit key removes that failure by
construction, and every arm inherits it.

\subsection{Two passes, and why}

Integration runs as two kernels over the same points.

\textbf{Allocate} walks the truncation band for each point, computes the block
coordinate of each voxel it touches, and ensures a block exists for it. This is
the irregular stage: threads probe for a data-dependent number of steps, race
for the same slots, and must observe one another's insertions.

\textbf{Update} walks the same band and accumulates into blocks that now exist.
This is the regular stage: a lookup that almost always hits on the first probe,
followed by five atomic additions per voxel.

Accumulation stores weighted \emph{sums} rather than running means. A running
mean requires reading the current value, combining, and writing back, which no
single atomic can make safe across five arrays; sums commute, so plain atomic
additions suffice and the arms need not agree on a locking scheme. The mean is
recovered at extraction.

Both passes apply the same occlusion test. Applying it in only one leaves blocks
that are allocated but never receive a contribution, which does not change the
extracted surface but does change the block count, and the block count is one of
the cross-arm correctness checks.

\subsection{The insert protocol}

Listing~\ref{lst:insert} is the CUDA C++ insert, and is the shape every arm
implements.

\begin{lstlisting}[caption={The insert protocol. Both the probe loop and the
publication window shape the comparison in Section~\ref{sec:attribution}.},
label={lst:insert}]
for (probe = 0; probe < size; ++probe) {
    entry = table[(slot + probe) & mask];
    key   = entry.key;                      // plain, cacheable load
    if (key == want) {                      // block already exists
        idx = entry.block_idx;
        while (idx < 0)                     // wait for the winner to publish
            idx = *(volatile int*)&entry.block_idx;
        return idx;
    }
    if (key == EMPTY) {
        prev = atomicCAS(&entry.key, EMPTY, want);
        if (prev == EMPTY) {                // we won the slot
            idx = atomicAdd(block_count, 1);
            block_coord[idx] = {bx, by, bz};
            __threadfence();                // coords visible before the index
            atomicExch(&entry.block_idx, idx);
            return idx;
        }
        // lost the race; fall through and re-examine this slot
    }
}
\end{lstlisting}

Two features of this protocol drive everything that follows.

\paragraph{The probe is unbounded.} It runs until it finds the key or an empty
slot, so its trip count is data-dependent and varies per lane. Section
\ref{sec:expressiveness} shows this is not expressible in Triton at any price.

\paragraph{Publication is two steps.} The key is published by a compare-exchange
and the index by a later store, so a reader can observe a key whose index is
still absent and must wait. The window exists because a block's \emph{identity}
is known before insertion while its \emph{location} is assigned by a counter and
known only after winning. Publishing both atomically would need the location
first; the location is only yours if the exchange succeeds. That circularity, and
two attempts to remove it, are discussed in Section~\ref{sec:attribution}.

\subsection{What is excluded, and why}

Surface extraction is not part of the comparison. It exists once, in CUDA C++,
and all three GPU arms call it, so measuring it three times compares nothing.
Making it a real comparison would mean writing marching tetrahedra twice more,
and that work would not extend the argument: it is a regular kernel, and the
comparison already contains one in the update stage.

Investigating the decision also corrected a number. The extraction stage had been
reported at 1.825\,ms and as several times the cost of the whole integrate path.
Splitting the call showed that 82\% of that measurement on one GPU and 97\% on the
other was a device-to-host transfer over PCIe links of differing width, not
compute. The kernel and its allocation are 0.319\,ms, comparable to integrate
rather than several times it, and the striking cross-device difference was the
motherboard.

\subsection{Why this workload}

The properties that make TSDF fusion awkward are exactly the ones absent from
the benchmarks these languages are usually compared on:

\begin{itemize}
  \item work per thread is data-dependent and not known at launch;
  \item threads in a warp diverge on both probe depth and outcome;
  \item correctness depends on compare-exchange and on one thread observing
        another's publication;
  \item the access pattern is a scatter through a hash, with no tile structure
        to exploit.
\end{itemize}

None of this is exotic. It is what any GPU hash table, sparse structure builder,
free-list allocator, or dynamic work queue looks like. The results in
Section~\ref{sec:results} should be read as applying to that class, not to
reconstruction specifically.

\section{Implementations}
\label{sec:impl}

% TODO: the five arms, the shared C ABI, the analytic correctness gate, and the
% rule that no code was reused between arms.

\section{Methodology}
\label{sec:method}

% TODO: batched launches, interleaved arm order, exclusivity guard, three passes
% and medians, the validity gates, and the dozen caught errors with the three
% that would have reversed a conclusion.

\section{Results}
\label{sec:results}

Table~\ref{tab:stage} is the paper in one table.

\begin{table}[h]
\centering
\begin{tabular}{llrr}
\toprule
Stage & Character & Rust / CUDA C++ & Triton / CUDA C++ \\
\midrule
allocate & irregular: probe, CAS, publication & 1.02--3.33 & 11.2--31.7 \\
update & regular: walk and accumulate & 0.96--1.18 & 1.1--2.5 \\
\bottomrule
\end{tabular}

\caption{Ratio to hand-written CUDA C++ by stage, across every recorded cell of
the sweep and both GPUs. The regular stage separates the languages by a small
factor; the irregular stage separates them by more than an order of magnitude.
The wide ends of both ranges are the smallest cell, where the kernel is short
enough that launch effects dominate; see Section~\ref{sec:results}.}
\label{tab:stage}
\end{table}

\begin{table}[h]
\centering
\small
\begin{tabular}{llrrrrr}
\toprule
Workload & GPU & CUDA C++ & \multicolumn{2}{c}{Rust} & \multicolumn{2}{c}{Triton} \\
\cmidrule(lr){4-5}\cmidrule(lr){6-7}
 & & ms & ms & $\times$ & ms & $\times$ \\
\midrule
RetroOffice P000, warm & RTX 5070 Ti & 0.388 & 0.396 & 1.02 & 1.134 & 2.92 \\
RetroOffice P000, warm & RTX 5060 & 0.763 & 0.768 & 1.01 & 2.229 & 2.92 \\
RetroOffice P000, cold & RTX 5070 Ti & 0.472 & 0.487 & 1.03 & 1.214 & 2.57 \\
RetroOffice P000, cold & RTX 5060 & 0.904 & 0.915 & 1.01 & 2.334 & 2.58 \\
AmericanDiner P000, warm & RTX 5070 Ti & 0.400 & 0.418 & 1.04 & 1.143 & 2.86 \\
AmericanDiner P000, warm & RTX 5060 & 0.780 & 0.796 & 1.02 & 2.103 & 2.70 \\
Plane, 320k points & RTX 5070 Ti & 0.363 & 0.380 & 1.05 & 0.876 & 2.41 \\
Plane, 320k points & RTX 5060 & 0.673 & 0.686 & 1.02 & 1.665 & 2.47 \\
Sphere, 1.28M points & RTX 5070 Ti & 0.975 & 1.024 & 1.05 & 2.986 & 3.06 \\
Sphere, 1.28M points & RTX 5060 & 1.712 & 1.735 & 1.01 & 5.706 & 3.33 \\
Sphere, 320k points & RTX 5070 Ti & 0.278 & 0.308 & 1.11 & 0.828 & 2.98 \\
Sphere, 320k points & RTX 5060 & 0.511 & 0.539 & 1.05 & 1.526 & 2.99 \\
\bottomrule
\end{tabular}

\caption{Full integrate path, allocate plus update, median of three passes.
Real depth data (TartanAir) and synthetic scenes, on two GPUs differing
$2.33\times$ in streaming multiprocessor count.}
\label{tab:totals}
\end{table}

% TODO: prose for the results, the workload sweep axes, and the device axis.

\section{Attribution}
\label{sec:attribution}

% TODO: Triton's two mechanisms; the scratch-region scaling failure; Rust's L1
% residency finding and the eleven eliminated hypotheses.

\section{What each language cannot say}
\label{sec:expressiveness}

% TODO: per-lane early exit, mask on atomic_cas, unbounded probe inexpressible,
% scoped atomic load versus caching, and the compiler defect found upstream.

\section{Threats to validity}
\label{sec:threats}

% TODO: one architecture (sm_120), the load-factor cap at 0.283, the
% unattributed residual, one real sequence.

\section{Related work}
\label{sec:related}

% TODO: Triton (Tillet et al., MAPL 2019); GPU language comparisons; TSDF
% fusion (KinectFusion, Voxel Hashing); Rust on GPU.

\section{Conclusion}
\label{sec:conclusion}

% TODO.

\end{document}
